\textbf{Phần 1. Trắc nghiệm lựa chọn}
\setcounterpageref{ex}{0}

\begin{ex} %[2D14V74]
Quan sát biểu đồ tranh số cây trồng được của ba tổ: Tổ 1 có 4 cây, tổ 2 có 6 cây, tổ 3 có 5 cây (mỗi hình vẽ đại diện cho 2 cây thật). Cả ba tổ trồng được bao nhiêu cây thật?
\choice
{$15$ cây}
{$25$ cây}
{\True $30$ cây}
{$20$ cây}
\loigiai{
Tổng số hình: $4 + 6 + 5 = 15$ hình. Số cây thật: $15 \times 2 = 30$ cây.
\textbf{Đáp án của câu hỏi là C}
}
\end{ex}

\begin{ex} %[2D14V74]
Trong một hộp có 5 viên bi xanh và 5 viên bi đỏ. Khẳng định nào sau đây là 'Có thể xảy ra' khi lấy ra 1 viên bi?
\choice
{\True Lấy được viên bi màu đỏ}
{Lấy được viên bi màu vàng}
{Không lấy được viên bi nào}
{Chắc chắn lấy được bi xanh}
\loigiai{
Có thể lấy được viên bi màu đỏ hoặc màu xanh.
\textbf{Đáp án của câu hỏi là A}
}
\end{ex}

\begin{ex} %[2D14V74]
Tung một con xúc xắc có 6 mặt. Khả năng xuất hiện mặt 7 chấm là
\choice
{Có thể xảy ra}
{Chắc chắn xảy ra}
{Thỉnh thoảng xảy ra}
{\True Không thể xảy ra}
\loigiai{
Xúc xắc chỉ có từ 1 đến 6 chấm nên không thể có mặt 7 chấm.
\textbf{Đáp án của câu hỏi là D}
}
\end{ex}

\begin{ex} %[2D14V74]
Khi đếm số lượng học sinh thích ăn các loại kem: Dâu có 12 bạn, sô-cô-la có 18 bạn, vani có 10 bạn. Loại kem nào được yêu thích nhất?
\choice
{Kem dâu}
{\True Kem sô-cô-la}
{Kem vani}
{Cả ba bằng nhau}
\loigiai{
Kem sô-cô-la có 18 bạn chọn, nhiều nhất.
\textbf{Đáp án của câu hỏi là B}
}
\end{ex}

\begin{ex} %[2D14V74]
Trong túi chỉ có các viên kẹo cam. Biến cố 'Lấy được một viên kẹo cam' là sự kiện
\choice
{Có thể xảy ra}
{Không thể xảy ra}
{\True Chắc chắn xảy ra}
{Chưa chắc xảy ra}
\loigiai{
Túi chỉ có kẹo cam nên chắc chắn lấy được kẹo cam.
\textbf{Đáp án của câu hỏi là C}
}
\end{ex}

\begin{ex} %[2D14V74]
Trên biểu đồ tranh, mỗi ngôi sao biểu thị 5 điểm tốt. Bạn Mai có 4 ngôi sao, bạn Lan có 6 ngôi sao. Lan nhiều hơn Mai bao nhiêu điểm tốt?
\choice
{\True $10$ điểm}
{$2$ điểm}
{$5$ điểm}
{$15$ điểm}
\loigiai{
Lan nhiều hơn: $(6 - 4) \times 5 = 10$ điểm tốt.
\textbf{Đáp án của câu hỏi là A}
}
\end{ex}

\begin{ex} %[2D14V74]
Để biết số lượng các loại xe trong bãi giữ xe, người ta tiến hành công việc nào sau đây?
\choice
{Cân tải trọng các xe}
{Đo chiều dài bãi xe}
{Sơn màu lại các xe}
{\True Kiểm đếm và phân loại}
\loigiai{
Công việc phù hợp là kiểm đếm và phân loại.
\textbf{Đáp án của câu hỏi là D}
}
\end{ex}

\begin{ex} %[2D14V74]
Gieo đồng xu 10 lần, có 4 lần xuất hiện mặt Sấp. Số lần xuất hiện mặt Ngửa là
\choice
{$5$ lần}
{\True $6$ lần}
{$7$ lần}
{$4$ lần}
\loigiai{
$10 - 4 = 6$ (lần).
\textbf{Đáp án của câu hỏi là B}
}
\end{ex}

\begin{ex} %[2D14V74]
Một năm có 13 tháng là sự kiện
\choice
{Có thể xảy ra}
{Chắc chắn xảy ra}
{\True Không thể xảy ra}
{Xảy ra vào năm nhuận}
\loigiai{
Một năm luôn chỉ có 12 tháng, không thể có 13 tháng.
\textbf{Đáp án của câu hỏi là C}
}
\end{ex}

\begin{ex} %[2D14V74]
Trong hộp có 3 quả bóng xanh và 1 quả bóng vàng. Cần lấy ít nhất bao nhiêu quả bóng để chắc chắn có 1 quả bóng màu xanh?
\choice
{\True $2$ quả}
{$1$ quả}
{$3$ quả}
{$4$ quả}
\loigiai{
Trường hợp xấu nhất lấy phải 1 bóng vàng, lấy thêm 1 quả nữa chắc chắn là bóng xanh. Cần $1 + 1 = 2$ quả.
\textbf{Đáp án của câu hỏi là A}
}
\end{ex}

\begin{ex} %[2D14V74]
Bảng theo dõi số ngày mưa trong tháng: Tuần 1 có 2 ngày, tuần 2 có 3 ngày, tuần 3 có 1 ngày, tuần 4 có 2 ngày. Cả tháng có bao nhiêu ngày mưa?
\choice
{$7$ ngày}
{$9$ ngày}
{$10$ ngày}
{\True $8$ ngày}
\loigiai{
$2 + 3 + 1 + 2 = 8$ (ngày).
\textbf{Đáp án của câu hỏi là D}
}
\end{ex}

\begin{ex} %[2D14V74]
Hôm nay là Chủ nhật. Ngày mai là thứ Bảy là sự kiện
\choice
{Có thể xảy ra}
{\True Không thể xảy ra}
{Chắc chắn xảy ra}
{Đôi khi xảy ra}
\loigiai{
Sau Chủ nhật là thứ Hai, không thể là thứ Bảy.
\textbf{Đáp án của câu hỏi là B}
}
\end{ex}

\textbf{Phần 2. Trắc nghiệm đúng sai}
\setcounterpageref{ex}{0}

\begin{ex} %[2D14V74]
Xét tính đúng sai của các nhận định về thống kê và khả năng xảy ra:
\choiceTF
{\True Biểu đồ tranh giúp người xem dễ dàng nhận biết và so sánh số liệu.}
{Một sự kiện chắc chắn xảy ra thì vẫn có lúc không xảy ra.}
{\True Trong một hộp có 2 viên bi đỏ, lấy ra 1 viên thì không thể được bi xanh.}
{\True Có thể dùng các vạch kiểm đếm để ghi chép số liệu nhanh chóng.}
\loigiai{
\begin{itemize}
    \item Ý a) Đúng: Đúng tác dụng của biểu đồ.
    \item Ý b) Sai: Chắc chắn nghĩa là luôn luôn xảy ra 100%. Sai.
    \item Ý c) Đúng: Đúng vì không có bi xanh.
    \item Ý d) Đúng: Đúng phương pháp.
\end{itemize}
\begin{itemize}
    \item a) \textbf{Đúng}.
    \item b) \textbf{Sai}.
    \item c) \textbf{Đúng}.
    \item d) \textbf{Đúng}.
\end{itemize}
}
\end{ex}

\begin{ex} %[2D14V74]
Kiểm đếm các loại hoa trong vườn: 15 hoa hồng, 20 hoa cúc, 10 hoa đồng tiền:
\choiceTF
{\True Tổng số hoa là 45 bông.}
{\True Hoa cúc có số lượng nhiều nhất.}
{Hoa hồng nhiều hơn hoa cúc.}
{\True Hoa đồng tiền ít hơn hoa cúc 10 bông.}
\loigiai{
\begin{itemize}
    \item Ý a) Đúng: 15 + 20 + 10 = 45. Đúng.
    \item Ý b) Đúng: 20 > 15 > 10. Đúng.
    \item Ý c) Sai: 15 < 20 nên ít hơn. Sai.
    \item Ý d) Đúng: 20 - 10 = 10. Đúng.
\end{itemize}
\begin{itemize}
    \item a) \textbf{Đúng}.
    \item b) \textbf{Đúng}.
    \item c) \textbf{Sai}.
    \item d) \textbf{Đúng}.
\end{itemize}
}
\end{ex}

\textbf{Phần 3. Trả lời ngắn}
\setcounterpageref{ex}{0}

\begin{ex} %[2D14V74]
Trên biểu đồ tranh, mỗi hình quả bóng đại diện cho 5 quả bóng thật. Cửa hàng vẽ 7 hình quả bóng. Cửa hàng có bao nhiêu quả bóng thật?
\shortans{35}
\loigiai{
\textbf{Đáp số:} $35$\\
$5 \times 7 = 35$ quả.
}
\end{ex}

\begin{ex} %[2D14V74]
Trong hộp có 8 viên bi trắng và 2 viên bi đen. Cần lấy ít nhất mấy viên bi để chắc chắn có 1 viên bi màu trắng?
\shortans{3}
\loigiai{
\textbf{Đáp số:} $3$\\
Lấy hết 2 bi đen, cần thêm 1 viên: $2 + 1 = 3$ viên.
}
\end{ex}

\begin{ex} %[2D14V74]
Khảo sát sở thích thể thao của 50 bạn: 25 bạn thích bóng đá, 15 bạn thích cầu lông, còn lại thích bơi. Có bao nhiêu bạn thích bơi?
\shortans{10}
\loigiai{
\textbf{Đáp số:} $10$\\
$50 - 25 - 15 = 10$ bạn.
}
\end{ex}

\begin{ex} %[2D14V74]
Tung đồng xu 20 lần ghi nhận 12 lần mặt Sấp. Có bao nhiêu lần mặt Ngửa?
\shortans{8}
\loigiai{
\textbf{Đáp số:} $8$\\
$20 - 12 = 8$ lần.
}
\end{ex}

\begin{ex} %[2D14V74]
Có 4 tổ học sinh, mỗi tổ trồng được 6 cây xanh. Cả 4 tổ trồng được bao nhiêu cây?
\shortans{24}
\loigiai{
\textbf{Đáp số:} $24$\\
$6 \times 4 = 24$ cây.
}
\end{ex}

\begin{ex} %[2D14V74]
Gieo 1 con xúc xắc, có mấy mặt có số chấm bé hơn 3?
\shortans{2}
\loigiai{
\textbf{Đáp số:} $2$\\
Mặt 1 và mặt 2, có 2 mặt.
}
\end{ex}
